\documentclass[10pt]{article}
\usepackage[T1]{fontenc}
\usepackage[utf8]{inputenc}
\usepackage[french]{babel}
\usepackage{amsmath,amssymb,mathrsfs,mathpazo}
\usepackage{geometry}
\geometry{a4paper,margin=2cm}
\usepackage{graphicx,array,tabularx,tabularray,colortbl,multirow,multicol,booktabs,enumitem,pifont,ulem,fancybox,diagbox,cancel}
\usepackage{pstricks,pstricks-add,pst-plot,pst-tree,pst-node,pst-text,pst-func,pst-math,pst-bspline,pst-3dplot,pst-eucl,pst-solides3d}
\let\pscalculate\undefined
\usepackage{pst-bezier}
\usepackage{tikz,pgfplots,tkz-tab}
\usetikzlibrary{arrows,arrows.meta,patterns,calc,positioning,trees,decorations.markings,shapes.geometric,shapes.misc,shapes.arrows,intersections,decorations.pathmorphing,decorations.pathreplacing,backgrounds,shadows,babel}
\pgfplotsset{compat=1.18}
\usepgfplotslibrary{fillbetween}
\usepackage[np]{numprint}
\usepackage{listings,makecell,pas-tableur,twoopt}
\let\barre\undefined
\usepackage[pstricks]{tablvar}
\definecolor{darkgreen}{RGB}{0,100,0}
\providecommand{\red}{\color{red}}
\providecommand{\blue}{\color{blue}}
\providecommand{\green}{\color{green}}
\providecommand{\black}{\color{black}}
\providecommand{\white}{\color{white}}
\providecommand{\gray}{\color{gray}}
\providecommand{\cyan}{\color{cyan}}
\providecommand{\magenta}{\color{magenta}}
\providecommand{\yellow}{\color{yellow}}
\providecommand{\violet}{\color{violet}}
\providecommand{\orange}{\color{orange}}
\providecommand{\euro}{\texteuro}
\providecommand{\EUR}{\texteuro}
\providecommand{\ieme}{\textsuperscript{e}}
\providecommand{\degres}{\ensuremath{^\circ}}
\providecommand{\degree}{\ensuremath{^\circ}}
\providecommand{\vv}[1]{\overrightarrow{#1}}
\providecommand{\Coord}[1]{#1}
\providecommand{\decosix}{}
\providecommand{\starredbullet}{\ensuremath{\bullet}}
\pagestyle{empty}
\setlength{\parindent}{0pt}
\usepackage{adjustbox}\DeclareUnicodeCharacter{2236}{:}\DeclareUnicodeCharacter{2190}{\ensuremath{\leftarrow}}
\providecommand{\R}{\mathbb{R}}
\providecommand{\N}{\mathbb{N}}
\providecommand{\Z}{\mathbb{Z}}
\providecommand{\Q}{\mathbb{Q}}
\providecommand{\D}{\mathbb{D}}
\providecommand{\C}{\mathbb{C}}
\providecommand{\e}{{\mathrm{e}}}
\providecommand{\vect}[1]{\overrightarrow{#1}}
\providecommand{\vectt}[1]{\overrightarrow{\mathrm{#1}}}
\providecommand{\vvt}[1]{\overrightarrow{\mathrm{#1}}}
\providecommand{\np}[2][]{{#2~#1}}
\providecommand{\vv}[1]{{\overrightarrow{#1}}}
\renewcommand{\d}{{\mathrm{d}}}
\providecommand{\degres}{{^\circ}}
\providecommand{\degre}{{^\circ}}
\providecommand{\blue}{\color{blue}}
\providecommand{\red}{\color{red}}
\providecommand{\enskip}{\quad}
\providecommand{\llbracket}{\lbrack\!\lbrack}
\providecommand{\rrbracket}{\rbrack\!\rbrack}
\providecommand{\euro}{{\text{€}}}
\providecommand{\no}{{\text{n°}}}
\providecommand{\interval}[2]{{[#1;#2]}}
\providecommand{\intervFO}[2]{{[#1;#2[}}
\providecommand{\dotfill}{\ldots}

\providecommand{\ds}{\displaystyle}
\providecommand{\textding}[1]{#1}
\providecommand{\Oij}{\(\mathrm{O};\vec{\imath},\vec{\jmath}\)}
\providecommand{\Oijk}{\(\mathrm{O};\vec{\imath},\vec{\jmath},\vec{k}\)}
\providecommand{\Ouv}{\(\mathrm{O};\vec{u},\vec{v}\)}
\providecommand{\intervFF}[2]{\ensuremath{[#1;#2]}}
\renewcommand{\euro}{\ifmmode\text{\texteuro}\else\texteuro\fi}
\DeclareUnicodeCharacter{00B0}{\ifmmode^{\circ}\else\textdegree\fi}
\DeclareUnicodeCharacter{20AC}{\ifmmode\text{\texteuro}\else\texteuro\fi}
% Énoncés : Éducation nationale. Transcription et corrigés : APMEP, Denis Vergès et contributeurs. Découpage et classement : CoopMaths, Sébastien Lozano, Éric Elter et contributeurs. Adaptation de la présentation : Maths Explorer.
% CC BY-SA : https://coopmaths.fr/www/about/
% https://coopmaths.fr/alea/static/eam/2026/tex/eam_specifique_2026_06_asie_1.tex
\begin{document}
\noindent\textit{APMEP / CoopMaths. CC BY-SA. Présentation adaptée par Maths Explorer.}\par\medskip
Dans un lycée, des élèves de première souhaitent installer des vélos générateurs d'électricité dans la cafétéria pour recharger des téléphones portables.

	À l'aide des quatre documents ci-dessous, indiquer si chacune des affirmations suivantes est vraie ou fausse. \textbf{Chaque réponse devra être justifiée}.

	\bigskip

	\textbf{Affirmation 1 :} Le budget du lycée permet de financer cinq vélos générateurs d'électricité (vélo, chargeur et livraison).

	\bigskip

	\textbf{Affirmation 2 :} Il faut pédaler pendant 24 minutes à intensité modérée pour recharger un téléphone portable à l'aide d'un vélo générateur d'électricité.

	\bigskip

	\textbf{Affirmation 3 :} On choisit au hasard un élève du lycée.

	La probabilité que l'élève choisi soit un élève de première en sachant qu'il est favorable au projet est supérieure à un quart.

	\bigskip

	\textbf{Affirmation 4 :} Le projet des élèves de première va pouvoir être lancé.

	\subsection*{Documents}
	\textbf{Document 1 : Fiche d'un vélo générateur d'électricité}

	\begin{center}
	\fbox{\begin{minipage}{0.9\linewidth}
		\textbf{Tarifs pour l'installation de vélos générateurs d'électricité}
		\begin{itemize}[left=5mm]
			\item Coût d'un vélo sans chargeur : 300 €
			\item Coût du chargeur, pour chaque vélo : 45 €
			\item Coût de livraison par vélo : 60 €
		\end{itemize}

		\textbf{Énergie produite par un vélo générateur d'électricité}

		Lorsqu'une personne pédale à intensité modérée pendant une heure, le dispositif permet de récupérer 100 kJ (kilojoules).
	\end{minipage}}
	\end{center}

	\textbf{Document 2 : Énergie nécessaire pour charger quelques appareils}
	\begin{center}
		\begin{tblr}{width=0.9\linewidth,colspec={X[c,2.5]*4{X[c]}},hlines, vlines, column{1} = {bg=gray!30},rows={m}}
			Appareils électroniques & Ordinateur portable & Casque audio & Téléphone portable & Montre connectée \\
			Énergie moyenne nécessaire pour une charge complète
			 & 200 kJ & 18 kJ & 40 kJ & 5 kJ \\
		\end{tblr}
	\end{center}

	\needspace{4\baselineskip}
	\textbf{Document 3 : Résultats du sondage réalisé auprès des élèves du lycée sur les vélos générateurs d'électricité}
	\begin{center}
		\begin{tblr}{width=0.9\linewidth,colspec={*4{X[c,m]}},column{1}={bg=gray!30},row{1}={bg=gray!30},cell{1}{1}={bg=white},hline{1}={2-Z}{solid},hline{2-Z}={solid},vline{1}={2-Z}{solid},vline{2-Z}={solid}}
			          & Favorable & Défavorable & Total \\
			Seconde   & 260       & 170         & 430   \\
			Première  & 180       & 120         & 300   \\
			Terminale & 120       & 150         & 270   \\
			Total     & 560       & 440         & 1000
		\end{tblr}
	\end{center}

	\textbf{Document 4 : Message de la proviseure du lycée aux élèves porteurs du projet}

	\begin{center}
	\fbox{\begin{minipage}{0.9\linewidth}
	\textbf{Objet : Lancement du projet "vélos générateurs d'électricité" -- conditions à
	remplir}

	Chers élèves de première,

	Ce projet ne pourra être lancé que si les \textbf{deux critères suivants sont remplis} :

	\begin{itemize}
		\item \textbf{plus de la moitié des élèves du lycée} doivent se déclarer favorables à cette initiative ;

		\item \textbf{au moins $60\,\%$ des élèves de première} doivent également y être \textbf{favorables}.

		Dans le cas où le projet serait validé, \textbf{le budget du lycée sera fixé à \np[\euro]{2000}} pour financer l'achat des vélos et des chargeurs ainsi que la livraison des vélos.
	\end{itemize}

	Bien cordialement,

	La Proviseure
	\end{minipage}}
	\end{center}

\end{document}
